\begin{frame}
    On considère le problème
\begin{equation*}
    \begin{cases}
        \partial_tu(t,x) - \nu\Delta u(t,x) - \alpha\Delta \big(\partial_t u(t,x)\big) = 0 & (t,x) \in \mathbb R_+ \times \bbR^d  \\
        u(0,x) = u_0(x) & x \in \bbR^d
    \end{cases}
\end{equation*}
\end{frame}

\subsection{Donnée initiale dans \texorpdfstring{$\scrS'(\bbR^d)$}{S'(Rd)}}
\begin{frame}{Unicité}
\begin{itemize}
\item On considère $u_0 \in \scrS'(\bbR^d)$. On peut reformuler le 
problème précédent de la façon suivante : trouver $u \in 
\scrC^1(\bbR_+, \scrS'(\bbR^d))$ tel que
\begin{equation}\label{pb_RS4}
    \begin{cases}
        u'(t) - \nu\Delta u(t) - \alpha\Delta \big(u'(t)\big) = 0  \\
        u(0) = u_0
    \end{cases}
\end{equation}
\item On suppose qu'il existe une solution 
$u \in \scrC^1(\bbR_+,\scrS'(\bbR^d))$. On applique la 
transformée de Fourier au problème.% \eqref{pb_RS4}.
%\begin{equation*}
%    \calF\big(u'(t)\big)-\nu\calF\big(\Delta u(t)\big)-\alpha\calF\big(\Delta(u'(t))\big) = 0
%\end{equation*} 
%\item On montre que $\calF\big(u'(t)\big)=\hat{u}'(t)$ et  que $\calF\Big(\dfrac{\mathrm{d}}{\mathrm{d}t}\Delta u(t)\Big) 
%= -\|\xi\|^2\hat{u}'(t)$. 
\item On obtient le problème suivant
\begin{equation}\label{pb_RS_Fourier}
    \begin{cases}
        (1+\alpha\|\xi\|^2)\hat{u}'(t)+\nu\|\xi\|^2\hat{u}(t) = 0  \\
        \hat{u}(0) = \hat{u}_0
    \end{cases}
\end{equation}
\end{itemize}
\end{frame}

%\begin{frame}{Unicité}
%    \begin{itemize}
%\item On suppose qu'il existe une solution 
%$u \in \scrC^1(\bbR_+,\scrS'(\bbR^d))$. On applique la 
%transformée de Fourier au problème \eqref{pb_RS4}.
%\begin{equation*}
%    \calF\big(u'(t)\big)-\nu\calF\big(\Delta u(t)\big)-\alpha\calF\big(\Delta(u'(t))\big) = 0
%\end{equation*} 
%\item On montre que $\calF\big(u'(t)\big)=\hat{u}'(t)$ et  que $\calF\Big(\dfrac{\mathrm{d}}{\mathrm{d}t}\Delta u(t)\Big) 
%= -\|\xi\|^2\hat{u}'(t)$. 
%\item Finalement, on obtient le problème suivant
%\begin{equation}\label{pb_RS_Fourier}
%    \begin{cases}
%        (1+\alpha\|\xi\|^2)\hat{u}'(t)+\nu\|\xi\|^2\hat{u}(t) = 0  \\
%        \hat{u}(0) = \hat{u}_0
%    \end{cases}
%\end{equation}
%\end{itemize}
%\end{frame}

\begin{frame}{Unicité}
\begin{itemize}
    \item Donc 
\begin{equation*}
    \hat{u}(t)(\xi) = \exp\bigg(\frac{-\nu \|\xi\|^2}{1+\alpha\|\xi\|^2}t\bigg)\hat{u}_0(\xi)
\end{equation*} 
\item On pose alors 
\begin{equation}\label{sol_Fourier}
    \forall (t,x) \in \bbR_+\times \bbR^d, \ u(t,\xi) = \calF^{-1}\bigg(\exp\bigg(\frac{-\nu \|\xi\|^2}{1+\alpha\|\xi\|^2}t\bigg)\hat{u}_0(\xi)\bigg)(x)
\end{equation}
Ce qui démontre l'unicité de la solution de \eqref{pb_RS4}.
\end{itemize}
\end{frame}

\begin{frame}{Existence}
\begin{itemize}
\item On va maintenant montrer que $u$ définie 
par \eqref{sol_Fourier} est solution de classe $\scrC^1$ de 
\eqref{pb_RS4}. 
\item On pose 
\begin{equation*}
    \forall (t,\xi) \in \bbR_+ \times \bbR^d, \ \hat{\calK}(t,\xi) := \exp\Big(\dfrac{-\nu \|\xi\|^2}{1+\alpha\|\xi\|^2}t\Big)
\end{equation*} 
\item Le fait que $\hat{\calK}(t,\cdot)$ soit un multiplicateur de 
$\scrS'(\bbR^d)$ nous permet de montrer que $\hat{u}$ est 
dérivable et que $\hat{u}'$ est continue.
    \item On vérifie aisément que 
\begin{equation*}
    \hat{u}(0)=\hat{u}_0
\end{equation*}
et que $\hat{u}(t)$ est solution de \eqref{pb_RS_Fourier} pour 
tout $t \in \bbR_+$. 
\item En utilisant les propriétés de $\calF$, on en déduit que \eqref{sol_Fourier} est 
solution de \eqref{pb_RS4} pour tout $t \in \bbR_+$. 
\item Enfin, en 
utilisant la formule d'inversion de Fourier et le lien entre 
convolution et transformée de Fourier, on peut réécrire l'expression de la 
solution
\begin{eqnarray}\label{sol_S'}
     u(t,x) = \frac{1}{(2\pi)^{d/2}}K(t,x)*u_0(x)
\end{eqnarray}
\end{itemize}
\end{frame}

%\begin{frame}{Existence}
%    \begin{itemize}
%    \item On vérifie aisément que 
%\begin{equation*}
%    \hat{u}(0)=\hat{u}_0
%\end{equation*}
%et que $\hat{u}(t)$ est solution de \eqref{pb_RS_Fourier} pour 
%tout $t \in \bbR_+$. 
%\item En utilisant les propriétés de $\calF$, on en déduit que \eqref{sol_Fourier} est 
%solution de \eqref{pb_RS4} pour tout $t \in \bbR_+$. 
%\item Enfin, en 
%utilisant la formule d'inversion de Fourier et le lien entre 
%convolution et transformée de Fourier, on peut réécrire l'expression de la 
%solution
%\begin{eqnarray}\label{sol_S'}
%     u(t,x) = \frac{1}{(2\pi)^{d/2}}K(t,x)*u_0(x)
%\end{eqnarray}
%\end{itemize}
%\end{frame}

\subsection{Donnée initiale dans \texorpdfstring{$H^s(\bbR^d)$}{Hs(Rd)}}

\begin{frame}
\begin{itemize}
\item On considère maintenant $u_0 \in H^s(\bbR^d)$. Comme 
$H^s(\bbR^d) \subset \scrS'(\bbR^d)$, on peut utiliser les 
résultats de la partie précédente qui nous assure de l'existence 
d'une solution dans $\scrS'(\bbR^d)$. 
\item Il suffit donc de montrer 
que la solution définie par \eqref{sol_S'} est dans 
$\scrC^1(\bbR_+,H^s(\bbR^d))$ ce qui nous donnera l'unicité de 
la solution dans $H^s(\bbR^d)$.
    \item Comme $\calF$ réalise une isométrie entre $H^s(\bbR^d)$ et 
$L^2_s(\bbR^d)$, on a
\begin{eqnarray*}
    u \in \scrC^1(\bbR_+,H^s(\bbR^d)) &\iff& \hat{u} \in \scrC^1(\bbR_+, L^2_s(\bbR^d))
\end{eqnarray*} 
\item Comme $u_0 \in H^s(\bbR^d)$, $\hat{u}_0 \in L_s^2(\bbR^d)$ donc pour presque tout $\xi \in \bbR^d$, $t \mapsto \hat{u}(t,\xi)$ est continue. On a également
\begin{equation*}
    \sup_{t \in \bbR_+}|\hat{u}(t,\cdot)| \leq |\hat{\calK}(t,\cdot)||\hat{u}_0| \leq |\hat{u}_0| \in L^2_s(\bbR^d)
\end{equation*} 
\item Enfin, 
\begin{equation*}
    \sup_{t \in \bbR_+} |\hat{u}'(t)| \leq \frac{\nu\|\xi\|^2}{1+\alpha\|\xi\|^2}|\hat{\calK}(t,\cdot)||\hat{u}_0| \leq \frac{\nu}{\alpha}|\hat{u}_0| \in L^2_s(\bbR^d)
\end{equation*}
Ce qui nous assure par théorème de convergence dominée que $\hat{u} \in \scrC^1(\bbR_+,L^2_s(\bbR^d))$.
\item Donc \eqref{sol_S'} avec $u_0 \in H^s(\bbR^d)$ est solution dans $H^s(\bbR^d)$.
\end{itemize}
\end{frame}

%\begin{frame}
%    \begin{itemize}
%    \item Comme $\calF$ réalise une isométrie entre $H^s(\bbR^d)$ et 
%$L^2_s(\bbR^d)$, on a
%\begin{eqnarray*}
%    u \in \scrC^1(\bbR_+,H^s(\bbR^d)) &\iff& \hat{u} \in \scrC^1(\bbR_+, L^2_s(\bbR^d))
%\end{eqnarray*} 
%\item Comme $u_0 \in H^s(\bbR^d)$, $\hat{u}_0 \in L_s^2(\bbR^d)$ donc pour presque tout $\xi \in \bbR^d$, $t \mapsto \hat{u}(t,\xi)$ est continue. On a également
%\begin{equation*}
%    \sup_{t \in \bbR_+}|\hat{u}(t,\cdot)| \leq |\hat{\calK}(t,\cdot)||\hat{u}_0| \leq |\hat{u}_0| \in L^2_s(\bbR^d)
%\end{equation*} 
%\item Enfin, 
%\begin{equation*}
%    \sup_{t \in \bbR_+} |\hat{u}'(t)| \leq \frac{\nu\|\xi\|^2}{1+\alpha\|\xi\|^2}|\hat{\calK}(t,\cdot)||\hat{u}_0| \leq \frac{\nu}{\alpha}|\hat{u}_0| \in L^2_s(\bbR^d)
%\end{equation*}
%Ce qui nous assure par théorème de convergence dominée que $\hat{u} \in \scrC^1(\bbR_+,L^2_s(\bbR^d))$.
%\item Donc \eqref{sol_S'} avec $u_0 \in H^s(\bbR^d)$ est solution dans $H^s(\bbR^d)$.
%\end{itemize}
%\end{frame}

\subsection{Comportement de la solution et comparaison avec l'équation de la chaleur}

\begin{frame}
    On rappelle l'équation de la chaleur
\begin{equation*}
    \begin{cases}
        \partial_tv(t,x)-\nu\Delta v(t,x)=0 & (t,x) \in \bbR_+ \times \bbR^d \\
        v(0,x) = v_0(x) & x \in \bbR^d
    \end{cases}
\end{equation*}
dont la solution est donnée par 
\begin{equation*}
    \forall (t,x) \in \bbR_+ \times \bbR^d, \ v(t,x) = \calF^{-1}\big(\exp(-\nu t\|\xi\|^2\big)\hat{v}_0(\xi)\big)(x)
\end{equation*}
\end{frame}

\begin{frame}
\begin{itemize}
    \item Dans un premier temps, on raisonne à $t \in \bbR_+$ fixé.
    \item A hautes fréquences, on remarque qu'on a l'équivalent suivant
    \begin{eqnarray*}
        \hat{\calK}(t,\xi) &\underset{\|\xi\|\to+\infty}{\sim}& \exp\Big(-\frac{\nu}{\alpha}t\Big)
    \end{eqnarray*}
    Donc le noyau $\hat{\calK}$ agit comme une multiplication par un scalaire. Donc
    \begin{eqnarray*}
        \hat{u}(t,\xi) &\underset{\|\xi\|\to+\infty}{\sim}& \exp\Big(-\frac{\nu}{\alpha}t\Big)\hat{u}_0(\xi)
    \end{eqnarray*}
        \item A basses fréquences, on peut effectuer un développement limité qui nous amène à
    \begin{eqnarray*}
        \hat{\calK}(t,\xi) &\underset{\|\xi\|\to 0}{\sim}& \exp(-\nu\|\xi\|^2t)
    \end{eqnarray*}
    Donc on a 
    \begin{eqnarray*}
        \hat{u}(t,\xi) &\underset{\|\xi\|\to 0}{\sim}& \exp(-\nu\|\xi\|^2t)\hat{u}_0(\xi)
    \end{eqnarray*}
\end{itemize}
\end{frame}

\begin{frame}
    \begin{itemize}
    \item Maintenant, si on regarde quand $t \to +\infty$ alors à hautes fréquences, comme on a l'équivalent
\begin{eqnarray*}
    \hat{\calK}(t,\xi) &\underset{\|\xi\|\to+\infty}{\sim}& \exp\Big(-\frac{\nu}{\alpha}t\Big)
\end{eqnarray*}
le noyau $\hat{\calK}$ va décroître de manière exponentielle. 
\item Tandis qu'à basses fréquences, 
\begin{eqnarray*}
    \hat{\calK}(t,\xi) &\underset{\|\xi\|\to 0}{\sim}& \exp(-\nu\|\xi\|^2t)
\end{eqnarray*}
donc le noyau $\hat{\calK}$ décroît plus lentement.
    \item La dynamique de la solution est donc dominée par les basses fréquences. On en déduit l'équivalent de la solution en temps long
\begin{eqnarray*}
    \hat{u}(t,\xi) &\underset{t \to +\infty}{\sim}& \exp(-\nu\|\xi\|^2t)\hat{u}_0(\xi)
\end{eqnarray*} 
\item Donc en temps long, la solution de l'équation de Rayleigh-Stokes se comporte comme la solution de l'équation de la chaleur.
\end{itemize}
\end{frame}